\section{Resultados principales}\label{sec:results}

\subsection{La reasignación del empleo poco calificado}

El resultado central del artículo es que la reforma de 2022 desplazó la composición del
empleo poco calificado desde el trabajo asalariado formal hacia el trabajo independiente,
sin elevar la remuneración de los trabajadores a los que estaba dirigida. Presento primero
la evidencia sobre la composición del empleo, luego la evidencia de primera etapa sobre
dónde el piso es vinculante y, en tercer lugar, la evidencia sobre los salarios.

El efecto estimado sobre la probabilidad de empleo formal entre los ocupados del sector
privado es $-0.046$ (columna 1 del Cuadro~\ref{tab:main}): una caída de 4.6 puntos
porcentuales para los trabajadores poco calificados respecto de los calificados, alrededor
de una quinta parte de la tasa de formalidad de base del grupo, de 23 por ciento. La
semielasticidad implícita es grande (cerca de $-2$ respecto del aumento de 10.2 por
ciento, columna 4 del Cuadro~\ref{tab:main}), pero ello refleja lo reducida que es la base
formal de este grupo y no una respuesta inverosímilmente elástica. La estimación es precisa
con errores agrupados por departamento, y la prueba de tendencias previas no rechaza
tendencias paralelas ($p=0.15$); la Sección~\ref{sec:robustness} toma en serio la
preocupación por una deriva residual y muestra que las magnitudes ajustadas por esa deriva
siguen siendo económicamente grandes y coinciden con la estimación de diferencias en
diferencias sintéticas.
El estudio de eventos de la Figura~\ref{fig:event} muestra que la formalidad relativa de
los poco calificados cae inmediatamente después de la reforma, en alrededor de cuatro
puntos porcentuales en el transcurso de un año, y que la caída persiste. En sentido
inverso, la probabilidad de trabajo independiente aumenta en 3.6 puntos porcentuales. A
diferencia de la formalidad, la serie de trabajo independiente sí rechaza la prueba de
tendencias previas ($p=0.003$), por efecto de un único trimestre elevado a mediados de
2021, por lo que me apoyo en el margen de formalidad como el más limpio de los dos y trato
la estimación de trabajo independiente como su imagen especular, que la corrobora. Las
horas semanales en la ocupación principal aumentan moderadamente y la condición de tiempo
completo disminuye levemente, en línea con un desplazamiento hacia el trabajo por cuenta
propia, donde las horas están menos estandarizadas.

Dos piezas de evidencia independientes respaldan una lectura causal de esta reasignación.
Primero, su magnitud crece con la exposición: como muestra la
Sección~\ref{sec:heterogeneity}, tanto la caída de la formalidad como el aumento del
trabajo independiente son mayores en los departamentos donde el nuevo piso corta más hondo
en la distribución salarial previa a la reforma. Segundo, un diseño de exposición continua
en la línea de \citet{card_1992}, que descarta por completo el contraste por educación y
relaciona los resultados con el índice de Kaitz regional estandarizado previo a la
reforma, arroja el mismo panorama: una incidencia (\emph{bite}) regional una desviación
estándar mayor se asocia con una caída de la formalidad 1.8 puntos porcentuales mayor y con
un aumento del trabajo independiente 1.0 punto mayor después de la reforma, estimaciones
que resisten la inferencia por aleatorización basada en el diseño a nivel de departamento
(Sección~\ref{sec:robustness}). En conjunto, estos resultados describen una reasignación
de los trabajadores poco calificados desde los empleos asalariados formales hacia el
trabajo por cuenta propia, el margen clásico de informalización que destacan
\citet{jaramillo_lopez_2006} para el Perú y \citet{ham_2018_honduras} para Honduras.

\begin{table}[!tbp]\centering
\caption{Estimaciones de diferencias en diferencias de la reforma del salario mínimo de 2022}\label{tab:main}
\begin{threeparttable}
\footnotesize
\setlength{\tabcolsep}{4pt}
\begin{tabular}{lcccc}
\toprule
Resultado & (1) DiD TWFE & (2) Doblemente robusto & (3) $p$ tend. previa & (4) Elasticidad \\
\midrule
Log salario real por hora & -0.021 & -0.056 & 0.001 & -0.20 \\
 & (0.012) & (0.137) & & \\
 & \multicolumn{2}{c}{\footnotesize $N=69\,534$, $\bar{y}=1.744$} & & \\
Log ingreso real mensual & -0.010 & -0.047 & 0.079 & -0.10 \\
 & (0.013) & (0.097) & & \\
 & \multicolumn{2}{c}{\footnotesize $N=70\,691$, $\bar{y}=6.901$} & & \\
Empleo & -0.025$^{***}$ & 0.020 & 0.000 & -0.35 \\
 & (0.006) & (0.037) & & \\
 & \multicolumn{2}{c}{\footnotesize $N=278\,064$, $\bar{y}=0.707$} & & \\
Participación laboral & -0.008$^{*}$ & -0.015 & -- & -0.10 \\
 & (0.004) & (0.034) & & \\
 & \multicolumn{2}{c}{\footnotesize $N=278\,064$, $\bar{y}=0.744$} & & \\
Empleo formal & -0.046$^{***}$ & 0.006 & 0.148 & -1.99 \\
 & (0.007) & (0.040) & & \\
 & \multicolumn{2}{c}{\footnotesize $N=181\,317$, $\bar{y}=0.228$} & & \\
Trab. independiente & 0.036$^{***}$ & 0.046 & 0.003 & 0.91 \\
 & (0.006) & (0.045) & & \\
 & \multicolumn{2}{c}{\footnotesize $N=181\,317$, $\bar{y}=0.389$} & & \\
Horas semanales & 0.556$^{**}$ & -0.424 & 0.034 & 0.13 \\
 & (0.249) & (0.872) & & \\
 & \multicolumn{2}{c}{\footnotesize $N=179\,111$, $\bar{y}=40.655$} & & \\
Gana menos que el SM & -0.006 & -0.007 & 0.128 & -0.16 \\
 & (0.007) & (0.039) & & \\
 & \multicolumn{2}{c}{\footnotesize $N=70\,691$, $\bar{y}=0.352$} & & \\
\bottomrule
\end{tabular}
\begin{tablenotes}\footnotesize
\item Notas: Cada fila es una regresión distinta sobre la ENAHO 2021Q1--2024Q4, que excluye el trimestre de transición 2022Q2, parcialmente tratado. La columna (1) reporta el coeficiente DiD de efectos fijos bidireccionales sobre $\mathrm{Low}\times\mathrm{Post}$ de la ecuación~\eqref{eq:twfe}, con efectos fijos de departamento y trimestre y controles (edad, edad al cuadrado, sexo, área urbana, estado civil, años de educación). La columna (2) reporta el estimador DiD doblemente robusto de \citet{santanna_zhao_2020}, colapsado a pre/post; sus errores estándar se agrupan por departamento usando la función de influencia del estimador. La columna (3) es el valor $p$ de una prueba conjunta de que las interacciones del estudio de eventos previas a la reforma son cero. La columna (4) convierte la estimación de la columna (1) en una elasticidad respecto del aumento de 10.2 por ciento del salario mínimo (para los resultados en niveles, en relación con la media base $\bar{y}$). Los resultados salariales corresponden a asalariados del sector privado (empleados y trabajadores del hogar, ingresos mensuales equivalentes); la formalidad, las horas y el trabajo independiente corresponden a trabajadores del sector privado. Errores estándar agrupados por departamento (25 grupos) entre paréntesis; todas las estimaciones están ponderadas con los factores de expansión de la ENAHO. $^{*}p<0.1$, $^{**}p<0.05$, $^{***}p<0.01$.
\end{tablenotes}
\end{threeparttable}\end{table}


\subsection{La primera etapa: donde se fiscaliza, el piso es vinculante}

¿Dónde es realmente vinculante el piso legal? Precisar esto primero ayuda a interpretar
los resultados sobre remuneraciones que siguen. El Cuadro~\ref{tab:firststage} y la
Figura~\ref{fig:bunching} examinan la distribución de ingresos de los asalariados privados
\emph{formales}, el segmento en el que se concentra la fiscalización, en ventanas
simétricas de tres trimestres alrededor de la reforma. Antes de la reforma, 14 por ciento
de los asalariados formales poco calificados percibe un ingreso dentro de un margen de 2.5
por ciento del piso anterior de 930 soles, un pico visible a la manera de
\citet{cengiz_2019} y \citet{engbom_moser_2022}. Después de la reforma, el pico se traslada
prácticamente por completo: la masa en 930 se reduce a menos de 2 por ciento y aparece un
pico casi idéntico, de 14 por ciento, en el nuevo piso de 1025 soles. En otras palabras,
los empleadores del sector formal cumplieron con el nuevo mínimo.

La respuesta de cumplimiento se divide exactamente a lo largo de la frontera de la
fiscalización. El Panel B del Cuadro~\ref{tab:firststage} estima las diferencias en
diferencias para un indicador de remuneración inferior a 1025 soles, manteniendo el umbral
fijo en todos los periodos. Entre los asalariados formales, la proporción por debajo del
nuevo piso cae en cinco puntos porcentuales para los poco calificados respecto de los
calificados (marginalmente significativo, $p=0.10$); entre los asalariados
\emph{informales}, en cambio, \emph{aumenta} en cuatro puntos porcentuales ($p=0.002$): la
remuneración informal se alejó aún más por debajo de un piso que en la práctica no se le
aplica. Ambas respuestas se compensan en la estimación conjunta con umbral fijo (última
fila del Panel B), y por eso tampoco cambia el indicador de remuneración inferior al mínimo
legal del Cuadro~\ref{tab:main}. Junto con los resultados de reasignación, esta evidencia
describe un ajuste coherente entre dos sectores: el piso es vinculante y se respeta dentro
del sector formal, el empleo formal del grupo afectado se contrae y el margen desplazado
reaparece como trabajo independiente y empleo informal, cuya remuneración el piso no
alcanza.

\begin{figure}[!t]
\centering
\includegraphics[width=0.8\textwidth]{fig10_bunching.pdf}
\caption{Traslado del pico en el piso salarial entre los asalariados privados formales poco
calificados. Histogramas ponderados de los ingresos nominales equivalentes mensuales
(intervalos de 50 soles) en ventanas simétricas de tres trimestres antes
(2021Q3--2022Q1) y después (2022Q3--2023Q1) de la reforma. Las líneas discontinuas marcan
el salario mínimo anterior (930) y el nuevo (1025).}
\label{fig:bunching}
\end{figure}

\begin{table}[!tbp]\centering
\caption{Primera etapa: migración del pico y cumplimiento del salario mínimo}\label{tab:firststage}
\begin{threeparttable}
\footnotesize
\begin{tabular}{lccc}
\toprule
  & Ventana previa & Ventana posterior & \\
\midrule
\multicolumn{4}{l}{\emph{Panel A: proporción de asalariados formales a $\pm$2.5\% de cada mínimo}} \\
Todos los asalariados formales & & & \\
\quad En el mínimo anterior (930) & 0.114 & 0.011 & \\
\quad En el nuevo mínimo (1025) & 0.052 & 0.117 & \\
Asalariados formales de baja educación & & & \\
\quad En el mínimo anterior (930) & 0.142 & 0.016 & \\
\quad En el nuevo mínimo (1025) & 0.064 & 0.140 & \\
\midrule
\multicolumn{4}{l}{\emph{Panel B: DiD sobre la proporción que gana menos de 1025 (umbral fijo)}} \\
 & DiD & (s.e.) & $p$ tend. previa \\ \midrule
Asalariados formales & -0.050 & (0.030) & 0.099 \\
Asalariados informales & 0.042$^{***}$ & (0.012) & 0.068 \\
Todos los asalariados & 0.025 & (0.021) & 0.007 \\
\bottomrule
\end{tabular}
\begin{tablenotes}\footnotesize
\item Notas: El panel A reporta la proporción ponderada de asalariados privados formales (remuneración mensual equivalente) a $\pm$2.5 por ciento de cada mínimo legal, en ventanas simétricas de tres trimestres antes (2021Q3--2022Q1) y después (2022Q3--2023Q1) de la reforma. El panel B reporta estimaciones DiD de $\mathrm{Low}\times\mathrm{Post}$ para un indicador de remuneración mensual equivalente menor a 1025 soles (umbral fijo en todos los períodos), según condición de formalidad, con errores estándar agrupados por departamento y el valor $p$ conjunto de tendencias previas en la última columna. Factores de expansión de la ENAHO en todos los casos. $^{*}p<0.1$, $^{**}p<0.05$, $^{***}p<0.01$.
\end{tablenotes}
\end{threeparttable}\end{table}


\subsection{La respuesta salarial}

Si los empleadores cubiertos cumplieron con el nuevo piso, ¿aumentó la remuneración del
grupo poco calificado en su conjunto? En promedio, no. El coeficiente de la interacción
entre el indicador de baja calificación y el indicador posterior a la reforma es $-0.021$
para el logaritmo del salario real por hora, con un error estándar de $0.012$: negativo,
marginalmente no significativo e incompatible con que la reforma haya elevado la
remuneración relativa. La estimación para el logaritmo de los ingresos mensuales es
igualmente nula. Soy deliberadamente cauto sobre cuánto puede acotar esta regresión, por
una razón transparente que se reporta en la columna (3) del Cuadro~\ref{tab:main}: la
prueba conjunta de tendencias previas para el salario por hora rechaza ($p=0.001$), debido a
un único pico positivo en 2021Q2 ($k=-4$) y no a una tendencia sistemática. Una vez que se
admiten desviaciones de las tendencias paralelas, el contraste por educación no puede
acotar con precisión el efecto salarial en ninguna dirección: bajo la restricción de
magnitudes relativas de \citet{rambachan_roth_2023}, el intervalo robusto con tendencias
paralelas exactas llega a $+2.9$ por ciento, y bajo una extrapolación lineal de la
tendencia previa (levemente decreciente) se admiten valores positivos aún mayores. Por ello
no sustento la conclusión de que la reforma no elevó los salarios en la regresión salarial
promedio. La sustento en las dos piezas de evidencia sobre remuneraciones cuyos
diagnósticos son limpios y que no dependen de la tendencia previa por educación.

Primero, la proporción de asalariados que percibe menos que el piso legal, que era de 35
por ciento entre los asalariados poco calificados en la línea de base, no cae después de la
reforma (la estimación es $-0.006$ con $p=0.45$, y sus tendencias previas son planas,
$p=0.13$). Si las empresas hubieran cumplido con el nuevo piso para el trabajador marginal,
esta proporción habría caído mecánicamente. Segundo, la distribución de ingresos de la
Figura~\ref{fig:dist} muestra una masa visible cerca del mínimo anterior antes de la
reforma, pero ningún traslado apreciable de esa masa hacia el nuevo piso después, lo que
refleja el incumplimiento generalizado documentado en la Sección~\ref{sec:data}; del mismo
modo, las estimaciones por cuantiles incondicionales de la Sección~\ref{sec:heterogeneity}
no muestran ganancias en ningún decil, ni siquiera en los inferiores, donde un piso que se
hace cumplir tendría que comprimir la distribución.

El estimador doblemente robusto de \citet{santanna_zhao_2020} de la columna (2) del
Cuadro~\ref{tab:main} merece una lectura franca y no una nota al pie. Este estimador
colapsa los dieciséis trimestres en una sola comparación antes y después de la reforma y
omite los efectos fijos de departamento, y con errores estándar agrupados por departamento
resulta impreciso para todos los resultados: la estimación salarial es negativa ($-0.056$,
error estándar $0.137$) y la de formalidad es nula ($+0.006$, error estándar $0.040$), en
lugar del $-0.046$ de efectos fijos bidireccionales. Leo esto como el costo conocido de
descartar la estructura temporal \citep{bertrand_2004}: colapsar a dos periodos y
veinticinco conglomerados deja poca variación identificadora, y los intervalos de
confianza resultantes contienen tanto el cero como efectos del doble de mi estimación de
base. Por esta razón trato la columna doblemente robusta como una verificación
complementaria conservadora y no como la estimación preferida, y la corroboración en la
que me apoyo proviene, en cambio, de diseños que conservan variación independiente: la
estimación de diferencias en diferencias sintéticas sobre celdas departamentales y el
diseño de exposición regional continua de la Sección~\ref{sec:robustness}.

Que la reforma no haya elevado los salarios en la parte baja de la distribución no
sorprende en este contexto: descontada la inflación, el valor real del piso aumentó poco, la
mayoría de los trabajadores poco calificados es informal o trabaja de manera independiente,
y la fiscalización es escasa, de modo que el mínimo legal restringe la remuneración solo de
una minoría. La ausencia de respuesta en las remuneraciones contrasta con el patrón de
``efecto faro'' que \citet{jimenez_2023} encuentra para la reforma de 2016. Retomo estos
mecanismos, y el contraste con reformas anteriores, en la Sección~\ref{sec:discussion}.


\subsection{El nivel de empleo}

El efecto estimado sobre la tasa de empleo total es de $-2.5$ puntos porcentuales, lo que
implicaría una elasticidad del empleo respecto del salario mínimo de alrededor de $-0.35$,
más del doble de las estimaciones más altas de la literatura peruana previa
\citep{cespedes_2005}. No doy a esta estimación una interpretación causal, por tres razones
documentadas en la Sección~\ref{sec:robustness}: la serie de empleo presenta una clara
tendencia previa diferencial, asociada a la recuperación tras la pandemia y concentrada a
inicios de 2021 ($p<0.001$ en la prueba de tendencias previas); reformas placebo ubicadas
en el periodo previo producen ``efectos'' de magnitud y significancia similares; y, a
diferencia de los márgenes de composición, el efecto sobre el empleo no crece con la
exposición regional; en todo caso, es menor en los departamentos de alta incidencia. Al
excluir 2021Q1, la estimación puntual se reduce aproximadamente a la mitad y no sobrevive a
los procedimientos de inferencia más conservadores. Por lo tanto, interpreto el resultado
sobre el nivel de empleo como reflejo de dinámicas de recuperación diferenciales y no del
salario mínimo, y concentro el peso interpretativo en el desplazamiento de la composición,
cuyos diagnósticos son limpios.

\subsection{Transiciones individuales: el margen de entrada}\label{sec:transitions}

Los resultados de composición describen stocks netos. Para identificar el flujo bruto que
los origina, complemento los cortes transversales repetidos con la ENAHO Panel 2020--2024,
que vuelve a entrevistar anualmente a las mismas viviendas y enlaza a las personas entre
años (Sección~\ref{sec:data}; los detalles de construcción están en el Apéndice~A). Apilo
los cuatro pares de años consecutivos, de 2020--21 a 2023--24, y clasifico cada par respecto
de la reforma según el mes de entrevista de 2022: pares previos a la reforma; pares de la
ventana de reforma, cuyo estado de origen se mide antes de mayo de 2022 y cuyo destino se
mide después; y pares posteriores a la reforma. Entre los trabajadores con un empleo
privado formal en el año base, estimo unas diferencias en diferencias apiladas para cada
estado de destino, con efectos fijos de par y de departamento, controles individuales,
ponderadores del panel que corrigen la atrición y errores estándar agrupados por
departamento. El Cuadro~\ref{tab:transitions} reporta las proporciones de destino y las
estimaciones.

Dos hallazgos precisan el mecanismo. Primero, los trabajadores en su puesto no fueron
expulsados. Respecto de los trabajadores formales de alta educación, la probabilidad de
que un trabajador formal de baja educación siga siendo formal \emph{aumenta} a lo largo de
la ventana de reforma (8.3 puntos porcentuales, $p=0.055$), y la transición hacia el empleo
asalariado informal cae ($-4.8$ puntos, $p=0.076$); las transiciones hacia el trabajo
independiente y hacia el no empleo no se mueven. Esto es exactamente lo que implica la
primera etapa: dentro del segmento fiscalizado, los trabajadores en su puesto recibieron el
aumento obligatorio y conservaron su empleo. Segundo, la puerta de entrada se cerró. Entre
quienes tenían un empleo informal en el año base, la probabilidad de que un trabajador de
baja educación pase \emph{a} un empleo privado formal cae en 2.1 puntos porcentuales
($p=0.011$) en los pares posteriores a la reforma, casi la mitad de la tasa de entrada de
base del grupo, de 4.5 por ciento. La caída del stock de formalidad documentada en la
Sección~\ref{sec:results} es, por lo tanto, un fenómeno del margen de las contrataciones:
la reforma no disolvió los vínculos formales existentes, sino que redujo el ritmo al que se
creaban otros nuevos para los trabajadores de baja educación. Esta lectura es coherente con
la incidencia de los efectos de composición entre los jóvenes, las mujeres y los perceptores
secundarios de ingreso, los trabajadores con mayor probabilidad de ser entrantes marginales
y no trabajadores protegidos en su puesto.

Tres salvedades acotan esta evidencia. El panel es anual, por lo que no puedo fechar las
transiciones dentro del año más allá de la división por mes de entrevista; los pares de
referencia previos a la reforma incluyen transiciones desde el año de la pandemia, 2020,
que tuvieron una rotación inusual; y la muestra con base formal es modesta (unas 7000
observaciones de pares), de modo que la estimación de permanencia es solo marginalmente
significativa. La respuesta de entrada aparece en los pares posteriores a la reforma y no
en la ventana inmediata de la reforma, lo que es coherente con flujos de contratación que se
ajustan con rezago. Por ello leo la evidencia del panel como una corroboración cualitativa
del mecanismo (permanencia de los trabajadores en su puesto más una menor entrada a la
formalidad) y no como magnitudes precisas.

\begin{table}[!tbp]\centering
\caption{Transiciones individuales en torno a la reforma (Panel ENAHO 2020--2024)}\label{tab:transitions}
\begin{threeparttable}
\footnotesize
\setlength{\tabcolsep}{5pt}
\begin{tabular}{lccc}
\toprule
\multicolumn{4}{l}{\emph{Panel A: proporciones de destino de los trabajadores privados formales en el año base}} \\
 & Baja educación & Alta educación & Baja $-$ Alta \\
\midrule
\multicolumn{4}{l}{Pares previos a la reforma} \\
\quad Sigue formal & 0.626 & 0.694 & -0.067 \\
\quad Asalariado informal & 0.151 & 0.116 & 0.035 \\
\quad Trabajo independiente & 0.128 & 0.067 & 0.061 \\
\quad Otro ocupado & 0.033 & 0.037 & -0.004 \\
\quad No ocupado & 0.062 & 0.086 & -0.024 \\
\multicolumn{4}{l}{Pares de la ventana de la reforma} \\
\quad Sigue formal & 0.717 & 0.697 & 0.020 \\
\quad Asalariado informal & 0.081 & 0.082 & -0.001 \\
\quad Trabajo independiente & 0.123 & 0.110 & 0.013 \\
\quad Otro ocupado & 0.029 & 0.048 & -0.019 \\
\quad No ocupado & 0.050 & 0.063 & -0.013 \\
\multicolumn{4}{l}{Pares posteriores a la reforma} \\
\quad Sigue formal & 0.589 & 0.674 & -0.085 \\
\quad Asalariado informal & 0.109 & 0.050 & 0.059 \\
\quad Trabajo independiente & 0.213 & 0.150 & 0.063 \\
\quad Otro ocupado & 0.026 & 0.052 & -0.025 \\
\quad No ocupado & 0.062 & 0.075 & -0.012 \\
\midrule
\multicolumn{4}{l}{\emph{Panel B: DiD apilado de transiciones (Low $\times$ período, EF de par y departamento)}} \\
 & Low$\times$Reforma & Low$\times$Post & Media previa a la reforma (baja) \\ \midrule
Sigue formal privado & 0.083$^{*}$ & -0.022 & 0.626 \\
 & (0.041) & (0.022) & \\
A asalariado informal & -0.048$^{*}$ & 0.020 & 0.151 \\
 & (0.026) & (0.014) & \\
A trabajo independiente & -0.005 & 0.009 & 0.128 \\
 & (0.017) & (0.011) & \\
A no ocupación & -0.010 & 0.011 & 0.062 \\
 & (0.012) & (0.013) & \\
Informal $\to$ formal (ingreso) & -0.006 & -0.021$^{**}$ & 0.045 \\
 & (0.005) & (0.008) & \\
\bottomrule
\end{tabular}
\begin{tablenotes}\footnotesize
\item Notas: Personas enlazadas del Panel ENAHO 2020--2024 (encuesta INEI 978), pares de años $t\to t{+}1$ para $t=2020,\dots,2023$, factores de expansión de panel \texttt{facpanel}. Los pares se clasifican respecto de la reforma del 1 de mayo de 2022 según el mes de entrevista de 2022: previos a la reforma (2020--21; 2021--22 entrevistados entre enero y abril de 2022), de la ventana de la reforma (2021--22 entrevistados desde mayo de 2022; 2022--23 entrevistados entre enero y abril de 2022, cuyo estado base es anterior a la reforma) y posteriores a la reforma (2022--23 entrevistados desde mayo de 2022; 2023--24). El panel A reporta las proporciones ponderadas de destino de los trabajadores con un empleo privado formal en el año base (7011 pares-observación). El panel B reporta coeficientes de regresiones apiladas de cada indicador de transición sobre Low$\times$Reform, Low$\times$Post y Low, con efectos fijos de par y departamento y con edad, edad al cuadrado y sexo como controles; la fila de ingreso se condiciona al empleo informal (asalariado o independiente) en el año base (33\,076 observaciones). Los estados de empleo replican las definiciones del texto principal; la informalidad usa el indicador oficial donde está publicado y la reconstrucción en los demás casos. Errores estándar agrupados por departamento del año base entre paréntesis. $^{*}p<0.1$, $^{**}p<0.05$, $^{***}p<0.01$.
\end{tablenotes}
\end{threeparttable}\end{table}


\begin{figure}[!t]
\centering
\includegraphics[width=0.98\textwidth]{fig5_event_study.pdf}
\caption{Estimaciones del estudio de eventos de la reforma de 2022. Cada panel grafica los
coeficientes $\beta_k$ de la interacción entre el indicador de baja calificación y los
indicadores trimestrales de la ecuación~\eqref{eq:event}, respecto del trimestre omitido
2022Q1 ($k=-1$). Se excluye el trimestre de transición 2022Q2 ($k=0$), parcialmente
tratado. Las bandas sombreadas son intervalos de confianza al noventa y cinco por ciento con
errores estándar agrupados por departamento.}
\label{fig:event}
\end{figure}

\begin{figure}[!t]
\centering
\includegraphics[width=0.72\textwidth]{fig4_wage_distribution.pdf}
\caption{Distribución de ingresos de los asalariados poco calificados, antes y después de
la reforma. Densidades kernel ponderadas de los ingresos nominales equivalentes mensuales
(empleados del sector privado y trabajadores del hogar), con el salario mínimo anterior
(930) y el nuevo (1025) marcados con líneas discontinuas.}
\label{fig:dist}
\end{figure}
